\chapter{Penutup}

\section{Kesimpulan}
Berdasarkan kegiatan MBKM di PT Vortex Buana Edumedia dan pembahasan pada bab-bab sebelumnya, dapat disimpulkan bahwa magang memberikan pengalaman langsung dalam pengembangan perangkat lunak pada lingkungan industri. Penulis berkontribusi pada pengembangan dokumentasi dan \textit{design library} Neon Interfaces serta fitur Import Nilai Neon Uji pada repositori Pendidikan. Kedua pekerjaan tersebut memiliki fokus berbeda, tetapi sama-sama menuntut ketelitian dalam memahami kebutuhan, merancang alur, mengembangkan fitur, menguji modul, dan berkomunikasi dengan tim.

Secara rinci, kesimpulan dari kegiatan ini adalah sebagai berikut:
\begin{enumerate}
    \item Mata kuliah Internship: Pengembangan Fitur dan Modul Proyek terlaksana melalui pengembangan dokumentasi komponen pada Neon Interfaces serta pengembangan fitur Import Nilai Neon Uji. Kegiatan tersebut memperkuat kemampuan menerjemahkan kebutuhan pengguna menjadi rancangan dan implementasi perangkat lunak yang relevan dengan kebutuhan perusahaan.
    \item Mata kuliah Internship: Pengujian Unit dan Modul Proyek terlaksana melalui penyusunan skenario uji, validasi fitur, debugging, dan pengecekan ulang pada modul yang dikembangkan. Pengujian membantu memastikan alur import nilai, pratinjau data, publikasi nilai, dan dokumentasi komponen memiliki perilaku yang konsisten.
    \item Mata kuliah Soft Skill: Kemampuan Berkomunikasi terlaksana melalui daily stand-up, diskusi design library, diskusi teknis fitur Import Nilai Neon Uji, training session intern, pelaporan progres, penyampaian kendala, dan penerimaan umpan balik. Kemampuan komunikasi terbukti penting untuk menjaga keselarasan antara pekerjaan teknis dan kebutuhan produk.
    \item Kegiatan magang memperlihatkan bahwa keberhasilan proyek perangkat lunak tidak hanya bergantung pada kemampuan implementasi, tetapi juga pada pemahaman domain, kualitas pengujian, dokumentasi yang baik, dan efektivitas komunikasi di dalam tim.
\end{enumerate}

Secara keseluruhan, kegiatan MBKM ini membantu menghubungkan pengetahuan akademik dengan praktik industri. Pengalaman tersebut menjadi bekal penting untuk memahami proses kerja profesional, standar kualitas perangkat lunak, dan tanggung jawab pengembang dalam menghasilkan sistem yang dapat digunakan oleh organisasi.

\section{Saran}
Saran yang dapat diberikan berdasarkan pelaksanaan kegiatan magang adalah sebagai berikut:
\begin{enumerate}
    \item Mahasiswa yang mengikuti kegiatan magang sebaiknya aktif mencatat kebutuhan, keputusan teknis, dan umpan balik selama proyek berlangsung agar proses pengembangan dan pengujian dapat ditelusuri dengan baik.
    \item Dalam pengembangan proyek perangkat lunak, pengujian sebaiknya dilakukan secara bertahap sejak awal implementasi. Cara ini membantu menemukan kesalahan lebih cepat dan mengurangi risiko perubahan yang merusak fitur lain.
    \item Dokumentasi komponen dan fitur sebaiknya dibuat seiring proses pengembangan. Dokumentasi yang baik membantu anggota tim lain memahami penggunaan komponen, alur fitur, dan batasan sistem.
    \item Komunikasi dengan pembimbing dan anggota tim perlu dilakukan secara rutin. Klarifikasi kebutuhan dan penyampaian progres yang jelas dapat membantu menjaga hasil pekerjaan tetap sesuai dengan tujuan proyek.
    \item Untuk pelaksanaan MBKM berikutnya, dokumentasi kegiatan dan keterkaitan dengan mata kuliah sebaiknya disusun sejak awal agar proses penyusunan laporan menjadi lebih terarah dan lengkap.
\end{enumerate}
